% Chapter 11 — what comes next: KG 7.1.0 (the entries' own images, Wikidata's facts), open gaps.
% Then the epilogue: the closing slide, to the right of the map, and the map once more.

\gmenter{next}
\note{Chapter 11: what is built and running now, and what is open.}

\gmgo{next-710}{7.1.0 and IMKG's questions}
\begin{frame}{7.1.0: each meme's image, and Wikidata's facts}
  \begin{columns}[t]
    \column{.37\textwidth}
      \gmpill{the meme's own image}\par\vskip2mm
      {\normalsize The vision model reads every entry's picture, not only its templates.}\par\vskip4mm
      \gmpill{Wikidata's facts}\par\vskip2mm
      {\normalsize Each linked item brings its facts: Kabosu \emph{is a} dog \dots}\par\vskip2mm
      \gmbig{\FStatements}{facts about \FStatementItems\ items}
    \column{.6\textwidth}
      {\normalsize\color{gmInk2}IMKG's four published questions, and its answers}\par\vskip3mm
      \begin{tikzpicture}[x=1cm, y=-1.05cm,
          q/.style={anchor=west, font=\normalsize\bfseries, text=gmInk, text width=4.0cm, align=left},
          a/.style={anchor=west, font=\normalsize, text=gmInk2, text width=4.1cm, align=left}]
        \node[q] at (0,0) {memes showing SpongeBob};   \node[a] at (4.3,0) {130};
        \node[q] at (0,1) {the most meme-able people}; \node[a] at (4.3,1) {Trump, Kyle Craven, Kanye West};
        \node[q] at (0,2) {memes based on films};      \node[a] at (4.3,2) {413};
        \node[q] at (0,3) {people by sex or gender};   \node[a] at (4.3,3) {male 10,333, female 2,798};
      \end{tikzpicture}\par\vskip3mm
      {\normalsize On 7.1.0: \textbf{the same queries, word for word} --- answered when it is published.}
  \end{columns}
  \note{
  \begin{itemize}
    \item Version 7.1.0 is built and tested; the image reading is running on the lab's GPU now
      (every entry's picture, about two days), and the graph will be published when it finishes.
    \item The meme's own image: IMKG used Google Vision on the meme's main picture; MemeAtlas reads
      it with the same vision model as the templates, and links what it sees from the meme itself.
    \item Wikidata facts: 753,854 statements of 32,494 linked items (1,183 properties), read from the
      same Wikidata dump as the lexicon. Now ``which memes show a human?'' is a question the graph
      can answer.
    \item The four use cases of IMKG's paper (Table 4), with IMKG's answers from its own notebook.
      They did not run on 7.0.0: three need Wikidata facts (is a human, is a film, sex or gender),
      one needs what the meme's image shows. 7.1.0 adds exactly that --- and because the vocabulary
      was kept, the queries run unchanged. [Update this slide when 7.1.0 is published:
      scripts/extract.py on the 7.1.0 build.]
  \end{itemize}}
\end{frame}

\gmgo{next-open}{Open questions}
\begin{frame}{Open questions}
  \begin{columns}[t]
    \column{.55\textwidth}
      \begin{tikzpicture}[x=1cm, y=-1.5cm, baseline=(current bounding box.north),
          ic/.style={circle, minimum size=1.0cm, fill=gm@col@next, text=white, font=\fontsize{14}{14}\selectfont, inner sep=0pt},
          tx/.style={anchor=west, text width=6.4cm, font=\normalsize, text=gmInk, align=left}]
        \node[ic] at (0,0) {\faQuestion}; \node[tx] at (.7,0) {\textbf{What Wikidata does not know}\\[-.4mm]{\small\color{gmInk2}Cheems is there; most meme characters are not}};
        \node[ic] at (0,1) {\faClone};    \node[tx] at (.7,1) {\textbf{One event, many pages}\\[-.4mm]{\small\color{gmInk2}two entries telling one happening, not yet joined}};
        \node[ic] at (0,2) {\faHistory};  \node[tx] at (.7,2) {\textbf{An old address}\\[-.4mm]{\small\color{gmInk2}no IRI says where an entry used to live}};
      \end{tikzpicture}
    \column{.41\textwidth}
      \begin{tikzpicture}[baseline=(current bounding box.north)]
        \node[fill=gmTint, rounded corners=2mm, inner sep=3.2mm, text width=4.6cm, align=left] {
          \gmwhy{Each one is written down}\par\vskip1mm\normalsize
          with what it costs and the fix, so the next person starts from there.};
      \end{tikzpicture}
  \end{columns}
  \note{
  \begin{itemize}
    \item Gap 11: names with no Wikidata item --- 22,614 in frame text, a fifth of the named people
      and characters in template images (Morty Smith, Springtrap \dots). Fixes range from a bigger
      lexicon to minting our own items.
    \item Cross-entry event identity: two pages telling the same happening produce two events.
    \item Gap 14, closed this week: one entry, one address. What remains open is the old IRIs:
      Doge's frame is now \texttt{/sensitive/memes/doge}, IMKG's was \texttt{/memes/doge}. Neo4j lists
      the old address on the frame; RDF says nothing (owl:sameAs would --- a modelling decision).
    \item The common thread: every weakness is documented with its cost and its fix.
  \end{itemize}}
\end{frame}

\gmepilogue
\gmgo{thanks}{Thank you}
{\setbeamertemplate{background}{}
\begin{frame}[plain]
  \begin{tikzpicture}[remember picture, overlay]
    \coordinate (c) at ($(current page.center)+(4.3cm,-.1cm)$);
    \foreach \a/\r/\col in {15/3.05/kWikidata, 48/2.55/kTemplate, 88/2.95/kFrame, 128/2.6/kEvent,
                            205/3.0/kConcept, 248/2.55/kWikidata, 292/2.95/kEvent, 328/2.5/kFrame} {
      \draw[gmBase, line width=.8pt] (c) -- ($(c)+(\a:\r)$);
      \fill[\col] ($(c)+(\a:\r)$) circle (.2);}
    \node[circle, minimum size=4.3cm, draw=kFrame, line width=3pt,
          path picture={\node at (path picture bounding box.center)
                        {\includegraphics[height=4.5cm]{figures/img/doge}};}] at (c) {};
    \node[anchor=south west, text width=7.9cm, font=\fontsize{20}{25}\selectfont\bfseries, text=gmInk, align=left]
      at ($(current page.west)+(11mm,-2mm)$) {\hyphenpenalty=10000 From a 2023 snapshot to a graph that renews itself every month.};
    \node[anchor=north west, font=\Large, text=gmInk2] at ($(current page.west)+(11.5mm,-6mm)$) {Thank you. Questions?};
  \end{tikzpicture}
  \note{
  \begin{itemize}
    \item One sentence to end: IMKG was a snapshot; MemeAtlas is a process --- it reads every page in
      depth, says how sure it is of every derived fact, keeps IMKG's language, and renews itself
      every month.
    \item Thank Riccardo and the lab. Questions.
  \end{itemize}}
\end{frame}}

\gmoverview{One meme,\\ eleven stops}
\note{The map again, every node visited. Leave it on screen during questions.}
